\documentclass[10pt,a4paper]{amsart}
\usepackage[margin=19mm]{geometry}
\usepackage{amsmath,amssymb,mathtools}
\usepackage[scr=rsfs,cal=euler]{mathalfa}
\usepackage{xfrac}
\usepackage[unicode,hidelinks,bookmarks=false]{hyperref}
\newcommand{\ie}{i.e.\ }
\newcommand{\cf}{cf.\ }
\newcommand{\resp}{respectively\ }
\newcommand{\Subsection}{\subsection}
\providecommand{\nobreakdash}{\nobreak\hbox{-}}
\setlength{\emergencystretch}{2em}
\multlinegap0pt
\usepackage{bm}
\usepackage{bbm}
\usepackage{dsfont}
\usepackage{xspace}
\usepackage{cancel}
\usepackage{mathtools}
\usepackage{amsthm}
\makeatletter
\@ifundefined{theorem}{\newtheorem{theorem}{Theorem}}{}
\@ifundefined{lemma}{\newtheorem{lemma}[theorem]{Lemma}}{}
\makeatother
\title[Horizontal Kakeya maximal operators in finite Heisenberg gro]{Horizontal Kakeya maximal operators in finite Heisenberg groups: Exact exponents and applications}
\author{Thang Pham, Andrea Pinamonti, Dung The Tran, Boqing Xue}
\date{Source: arXiv:2603.02111v2 (CC BY 4.0)}
\begin{document}
\maketitle

\noindent\emph{Source and licence.} Horizontal Kakeya maximal operators in finite Heisenberg groups: Exact exponents and applications. Version arXiv:2603.02111v2 (CC BY 4.0), distributed under the Creative Commons Attribution 4.0 International licence, as recorded in the arXiv licence metadata for this exact version and shown on the abstract page https://arxiv.org/abs/2603.02111v2. Full rights notice: licenses/arxiv-2603.02111-CC-BY-4.0.txt.\par\smallskip

\noindent\textit{Editorial scope.} Counted unit: the Theorem whose source label is \texttt{thm:planar-all-u} in the article (source0.tex lines 922--938, source label \texttt{thm:planar-all-u}), delivered together with the article's own complete original proof of it (source0.tex lines 941--1095, one \texttt{proof} environment of 155 lines). No mathematics is written, repaired or re-derived: statement and proof are the article's own text line for line. Required context retained uncounted: interpolation-inequality (lines 612--635); lem:planar-l2 (lines 735--742). Every definition, display and label of those lines is retained unchanged. Omitted from this package and not redistributed: 1--611, 636--734, 743--921, 939--940, 1096--3293. Nothing inside a retained excerpt is omitted or altered: every retained body is delivered exactly as the article's lines are written, so no omission receipt of any kind accompanies this package. Citations: the retained excerpts carry no citation command at all, so no bibliography entry of the article is transcribed and no reference is redistributed. Reference adaptation: none was needed and none was made; every reference inside the delivered unit resolves inside it, so no unresolved reference or citation is produced. Printed slips kept literal: no printed word, symbol, hypothesis or number of the counted statement or of its proof was altered. Layout and class: the article's own class is \texttt{article} with packages including tikz, geometry, amsmath, amssymb, amsthm, comment, enumitem, array, color, svg, parskip, hyperref, setspace, graphicx; the wrapper is the lane's pinned \texttt{amsart} editorial wrapper at 10pt on A4, which restates the theorem-like environments the retained text needs and adds the article's own macro definitions that the retained text needs (none), with extra wrapper packages (bm, bbm, dsfont, xspace, cancel, mathtools). The delivered numbering is the unit's own. Assets: no retained span contains a figure environment, a graphics-inclusion command, a PDF-page inclusion, a table or a TikZ picture; no image file is copied into this package, the package declares no asset dependency and no image file is redistributed with it. Licence: version arXiv:2603.02111v2 of this article is distributed under the Creative Commons Attribution 4.0 International licence, as recorded in the arXiv licence metadata for this exact version and shown on the abstract page https://arxiv.org/abs/2603.02111v2. A full rights notice is delivered at licenses/arxiv-2603.02111-CC-BY-4.0.txt. Characters: every retained excerpt is pure ASCII, so the delivered engine, XeLaTeX with its default Latin Modern text font, reports no missing character.
\begin{lemma}\label{interpolation-inequality}
Let $X$ and $Y$ be finite sets, and let
\[
T : \ell^{p_0}(X) \cap \ell^{p_1}(X)\longrightarrow \ell^{r_0}(Y) \cap \ell^{r_1}(Y)
\]
be a linear operator.
Let $1 \le p_0, p_1, r_0, r_1 \le \infty$.
Assume that the following estimates hold for all functions $f : X \to \mathbb{C}$:
\[
\|Tf\|_{\ell^{r_0}(Y)} \le C_0 \, \|f\|_{\ell^{p_0}(X)},
\qquad
\|Tf\|_{\ell^{r_1}(Y)} \le C_1 \, \|f\|_{\ell^{p_1}(X)}.
\]
Then, for any $0 \le \theta \le 1$, we have
\[
\|Tf\|_{\ell^{r}(Y)} \le C_0^{1-\theta} C_1^{\theta} \, \|f\|_{\ell^{p}(X)},
\]
where the exponents $p$ and $r$ are defined by
\[
\frac{1}{p} = \frac{1-\theta}{p_0} + \frac{\theta}{p_1},
\qquad
\frac{1}{r} = \frac{1-\theta}{r_0} + \frac{\theta}{r_1}.
\]
\end{lemma}

\begin{lemma}\label{lem:planar-l2}
For every function $f:\mathbb{F}_q^2\to\mathbb{C}$, one has
\[
\|\mathcal{M}_{2}f\|_{\ell^2(\mathbb{P}^{1}(\mathbb{F}_q))}
\le
\sqrt{2}\,q^{\frac{1}{2}}\,\|f\|_{\ell^2(\mathbb{F}_q^2)}.
\]
\end{lemma}

\begin{theorem}
\label{thm:planar-all-u}
For every $1\le u\le\infty$ and for all $G:\mathbb{F}_q^2\to[0,\infty)$, we have 
\[
\|\mathcal{M}_{2}G\|_{\ell^u(\mathbb{P}^1(\mathbb{F}_q))}
\le \sqrt{2}
\,q^{\tau(u)}\,\|G\|_{\ell^u(\mathbb{F}_q^2)}
\]
where
\[
\tau(u):=
\begin{cases}
\displaystyle \frac1u, & 1\le u\le 2,\\[0.4em]
\displaystyle 1-\frac1u, & 2\le u\le \infty.
\end{cases}
\]
\end{theorem}

 \begin{proof}
Let $G:\mathbb{F}_q^2\to[0,\infty)$ be given.
For each direction $[\mathbf{v}]\in\mathbb{P}^1(\mathbb{F}_q)$, choose an affine line
$\ell_{[\mathbf{v}]}\subset\mathbb{F}_q^2$ with direction $[\mathbf{v}]$ such that
\[
\mathcal{M}_{2}G([\mathbf{v}])=\sum_{\mathbf{x} \in \ell_{[\mathbf{v}]}} G(\mathbf{x}).
\]
Define the linear operator $T$ on functions $f:\mathbb{F}_q^2\to\mathbb{C}$ by
\[
(T f)([\mathbf{v}]):=\sum_{\mathbf{x} \in\ell_{[\mathbf{v}]}} f(\mathbf{x}),
\qquad [\mathbf{v}]\in\mathbb{P}^1(\mathbb{F}_q).
\]
Then
\begin{equation}\label{eq:TG-linearizes}
(T G)([\mathbf{v}])=\mathcal{M}_{2}G([\mathbf{v}])
\qquad\text{for all }[\mathbf{v}]\in\mathbb{P}^1(\mathbb{F}_q).
\end{equation}
Therefore,
\begin{equation}\label{eq:reduce-to-TG}
\|\mathcal{M}_{2}G\|_{\ell^u(\mathbb{P}^1(\mathbb{F}_q))}
=
\|T G\|_{\ell^u(\mathbb{P}^1(\mathbb{F}_q))}.
\end{equation}

We next record three endpoint bounds for $T$.


\medskip
\noindent\textbf{The $\ell^1$ bound.}
For every $f:\mathbb{F}_q^2\to\mathbb{C}$,
\begin{align}\label{eq:TG-l1}
\|T f\|_{\ell^1(\mathbb{P}^1(\mathbb{F}_q))}
&=
\sum_{[\mathbf{v}]\in\mathbb{P}^1(\mathbb{F}_q)}
\Bigl|\sum_{\mathbf{x} \in\ell_{[\mathbf{v}]}} f(\mathbf{x})\Bigr|
\le
\sum_{[\mathbf{v}]}\sum_{\mathbf{x} \in\ell_{[\mathbf{v}]}}|f(\mathbf{x})|
\le
(q+1)\sum_{\mathbf{x} \in\mathbb{F}_q^2}|f(\mathbf{x})| \notag\\
&=
(q+1)\,\|f\|_{\ell^1(\mathbb{F}_q^2)}.
\end{align}

\medskip
\noindent\textbf{The $\ell^\infty$ bound.}
For every $f:\mathbb{F}_q^2\to\mathbb{C}$ and every $[\mathbf{v}]\in \mathbb{P}^1(\mathbb{F}_q)$,
\[
|(T f)([\mathbf{v}])|
=
\Bigl|\sum_{\mathbf{x} \in\ell_{[\mathbf{v}]}} f(\mathbf{x})\Bigr|
\le
\sum_{\mathbf{x} \in\ell_{[\mathbf{v}]}}|f(\mathbf{x})|
\le
q\,\|f\|_{\ell^\infty(\mathbb{F}_q^2)},
\]
since each affine line in $\mathbb{F}_q^2$ has exactly $q$ points.
Taking the maximum over $[\mathbf{v}]\in \mathbb{P}^1(\mathbb{F}_q)$ implies
\begin{equation}\label{eq:TG-linfty}
\|T f\|_{\ell^\infty(\mathbb{P}^1(\mathbb{F}_q))}
\le
q\,\|f\|_{\ell^\infty(\mathbb{F}_q^2)}.
\end{equation}

\medskip
\noindent\textbf{The $\ell^2$ bound.}
The operator $T$ is obtained by selecting one affine line in each direction,
so the proof of Lemma~\ref{lem:planar-l2} applies verbatim to $T$.
Therefore, for every $f:\mathbb{F}_q^2\to\mathbb{C}$,
\begin{equation}\label{eq:TG-l2}
\|T f\|_{\ell^2(\mathbb{P}^1(\mathbb{F}_q))}
\le
\sqrt{2}\,q^{\frac12}\,\|f\|_{\ell^2(\mathbb{F}_q^2)}.
\end{equation}

We now interpolate the bounds to obtain the desired exponents.

\medskip
\noindent\textbf{Case 1: $1\le u\le 2$.}
Choose $\theta\in[0,1]$ so that
\[
\frac1u=(1-\theta)\cdot 1+\theta\cdot\frac12.
\]
Then $\theta=2\bigl(1-\frac1u\bigr)$.
Applying Lemma~\ref{interpolation-inequality} to the linear operator $T$,
using the endpoint estimates \eqref{eq:TG-l1} and \eqref{eq:TG-l2}, we obtain
{ 
\begin{align*}
\|Tf\|_{\ell^u(\mathbb{P}^1(\mathbb{F}_q))}
\le
(q+1)^{1-\theta}\bigl(\sqrt{2}\,q^{\frac12}\bigr)^{\theta}
\|f\|_{\ell^u(\mathbb{F}_q^2)} 
\le
\left(\frac43\right)^{1-\theta}
2^{\frac{\theta}{2}}
q^{1-\frac{\theta}{2}}
\|f\|_{\ell^u(\mathbb{F}_q^2)},
\end{align*}
where the last inequality follows from the assumption that $q\ge 3$.

Since $1\le u\le2$, we have $\theta\in[0,1]$. Moreover,
\[
\left(\frac43\right)^{1-\theta}2^{\frac{\theta}{2}}
=
\frac43\left(\frac{3\sqrt2}{4}\right)^\theta,
\]
which is an increasing function of $\theta$ because $\frac{3\sqrt2}{4}>1$.
Therefore,
\begin{align}
\label{eq:TG-u-1-2}
\|Tf\|_{\ell^u(\mathbb{P}^1(\mathbb{F}_q))}
\le
\left(\frac43\right)^{1-\theta} \, 
2^{\frac{\theta}{2}} \, 
q^{1-\frac{\theta}{2}} \, 
\|f\|_{\ell^u(\mathbb{F}_q^2)}
\le
\sqrt{2}\,
q^{1-\frac{\theta}{2}} \, 
\|f\|_{\ell^u(\mathbb{F}_q^2)}.
\end{align}
}
From here, we apply \eqref{eq:TG-u-1-2} with $f=G$ and use \eqref{eq:reduce-to-TG}.



\medskip
\noindent\textbf{Case 2: $2\le u\le \infty$.}
Choose $\theta\in[0,1]$ so that
\[
\frac1u=(1-\theta)\cdot\frac12+\theta\cdot 0.
\]
Then $\theta=1-\frac{2}{u}$.
Applying Lemma~\ref{interpolation-inequality} to the linear operator $T$,
using the endpoint estimates \eqref{eq:TG-l2} and \eqref{eq:TG-linfty}, we obtain 
\[
\|T f\|_{\ell^u(\mathbb{P}^1(\mathbb{F}_q))}
\le
\bigl(\sqrt{2}q^{\frac12}\bigr)^{1-\theta}\,q^{\theta}\,
\|f\|_{\ell^u(\mathbb{F}_q^2)}
\leq \sqrt{2}
\,q^{\frac{1-\theta}{2}+\theta}\,\|f\|_{\ell^u(\mathbb{F}_q^2)}.
\]
Since $\frac{1-\theta}{2}+\theta=1-\frac1u$, this implies
\[
\|T f\|_{\ell^u(\mathbb{P}^1(\mathbb{F}_q))}
\le \sqrt{2} \,q^{1-\frac1u}\,\|f\|_{\ell^u(\mathbb{F}_q^2)}.
\]

Combining the two cases gives
\[
\|T f\|_{\ell^u(\mathbb{P}^1(\mathbb{F}_q))}
\le \sqrt{2}  \,q^{\tau(u)}\,\|f\|_{\ell^u(\mathbb{F}_q^2)}.
\]
Taking $f=G$ and combining \eqref{eq:reduce-to-TG}, the theorem then follows.
\end{proof}

\end{document}
